\section{Introduction}

The PatchOps problem sits at the intersection of several active research areas rather than belonging to only one
of them. It is related to log-based root-cause analysis, automated program repair, repository-level software
engineering agents, CI benchmark construction, failure taxonomy studies, and the growing literature on
calibration, abstention, and overfitting in LLM-based software systems
\cite{logsage2025,llmaprsurvey2025,abstentionSurvey2026,aprtop2025}. A good related-work chapter for
PatchOps therefore needs to do more than list nearby papers. It must explain which parts of the prior literature
PatchOps inherits directly, which claims have already been established by others, and which gap still remains
meaningful enough to support a thesis contribution.

As of mid-2026, the field has moved beyond the stage where one can claim novelty simply by saying that an LLM
reads CI logs and proposes a fix. Industrial and academic systems already demonstrate that CI diagnosis and even
CI-validated repository repair are feasible tasks. The important question is no longer whether repair should be
attempted at all, but rather how such repair should be \emph{framed}, \emph{controlled}, and \emph{evaluated}.
This chapter reviews the most relevant prior work and shows why PatchOps is best positioned not as ``the first
CI repair agent,'' but as a safety-gated, abstention-aware, closed-loop repair system built on top of an
already established task.

\section{LLM-Based CI/CD Diagnosis and Remediation}

The closest diagnosis-oriented system to PatchOps is \textbf{LogSage} (ASE 2025), an industrially validated
framework for CI/CD failure detection and remediation. LogSage combines an offline knowledge-building phase
with an online diagnosis phase in which failed logs are filtered, expanded, and pruned before being passed to
LLM prompts supported by retrieval \cite{logsage2025}. Its core strength is not repository patch generation, but industrial-scale
root-cause analysis with high reported precision and a large deployment footprint. For PatchOps, LogSage is
important for two reasons.

First, it establishes that log preprocessing and focused evidence extraction matter. Raw CI logs are noisy,
redundant, and often too large to feed directly to an LLM. PatchOps inherits this lesson through its log
extraction stage and should be understood as part of the same design tradition: extract the critical signal
before asking a model to reason.

Second, LogSage sets a strong baseline in the diagnosis portion of the problem. It would be unrealistic and
unnecessary for PatchOps to compete with a year-long industrial deployment on the narrow question of root-cause
analysis precision. The more useful interpretation is that systems such as LogSage solve an important part of
the diagnosis problem, while leaving open the question of how diagnosis should connect to repository edits,
patch validation, and safety-sensitive action selection.

The most directly relevant repair-oriented work is the reference workflow introduced with
\textbf{CI-Repair-Bench} (arXiv, April 2026) \cite{cirepairbench2026}. Its pipeline follows a simple but consequential structure:
analyze the CI log, localize the likely fault, generate a patch, and then re-run the original workflow to
measure whether the fix succeeds under true CI-style validation. This benchmark is a turning point for the
PatchOps thesis because it already establishes repository-aware, CI-verified repair as a public task. It also
provides a concrete performance reference, with the strongest baseline still achieving only around one-fifth
Pass@1 overall. In other words, CI repair has been formalized, but it remains far from solved.

This matters for thesis positioning. PatchOps cannot honestly claim to be the first system to validate patches
through CI, nor can it claim that repository-aware patch repair itself is unexplored. What remains open is what
CI-Repair-Bench does not emphasize: when an agent should decline to patch, how it should behave on transient or
risky failures, and how to prevent superficially successful but unsafe or off-target edits from being treated as
wins. That reframing is central to PatchOps.

\section{Automated Program Repair and Build Repair}

PatchOps also belongs to the longer history of automated program repair (APR), especially the branch of APR
concerned with build and compilation failures. Earlier systems such as \textbf{HireBuild} and
\textbf{DeepDelta} addressed narrower but structurally similar tasks: they used mined fix patterns or learned
transformations to repair failing build scripts or compilation errors. These systems demonstrated that repeated
software failure patterns can often be corrected automatically, especially when the failure class is narrow and
the repair space is constrained \cite{hirebuild2018,deepdelta2019}.

Their relevance to PatchOps is historical and conceptual. They show that repair is not a brand-new ambition,
but they also expose the limitations of narrow repair regimes. Build-script repair, compilation repair, and
configuration repair work best when the failure type is predictable and the patch space is small. CI repair is
harder because CI failures are broader in scope: they may arise from code, tests, dependencies, workflow YAML,
environment setup, permissions, missing secrets, or flaky infrastructure. PatchOps therefore inherits the
repair ambition of pre-LLM systems, but it operates in a setting that is more heterogeneous and less
structurally constrained.

The broader LLM-APR field makes this contrast even clearer. Surveys of LLM-based APR show fast progress in
patch generation, but they also show that success depends heavily on task framing, benchmark contamination,
validation quality, and the ability to localize the right files before editing them. Repository-level agentic
benchmarks such as SWE-bench and its descendants show that well-structured pipelines can outperform overly
complex multi-agent designs \cite{llmaprsurvey2025,llmaprslr2024,swebench2023,agentless2025}. This is an important design lesson for PatchOps. A staged state-machine pipeline,
with explicit transitions from log evidence to localization, patch planning, generation, validation, and final
decision, is not merely an engineering convenience. It is a research-informed choice that aligns with evidence
that simpler, structured repair workflows are often easier to reproduce, analyze, and ablate than multi-agent
systems built for appearance rather than clarity.

\section{Why CI Repair Is Not the Same as Classical APR}

Although PatchOps is related to APR, CI repair should not be reduced to a variant of classical test-suite-based
program repair. Traditional APR usually assumes that a developer test suite exists, that one can execute it
repeatedly, and that patch correctness is approximated by whether tests pass. CI repair is more operational and
more fragile. It may involve infrastructure scripts, package installation, linter configuration, workflow
permissions, or environmental mismatches rather than source-code semantics alone. A failing GitHub Actions run
is therefore not simply a failed unit test. It is a signal produced by a workflow that may stop at the first
failure and may expose only a partial view of the system state.

This distinction becomes especially important when considering correctness. Classic APR literature has spent
years documenting the gap between a \emph{plausible} patch and a \emph{correct} patch
\cite{llmaprslr2024,aprtop2025,compass2026}. The same warning applies,
and perhaps applies even more strongly, to CI repair. A patch that makes CI green is not automatically a good
patch. It may silence a check, change the wrong file, overfit to the immediate failure, or exploit a weakness
in the workflow itself. Because CI validation is a partial oracle, green should be interpreted as necessary but
not sufficient evidence. This is one of the strongest motivations for PatchOps' safety gate, patch-plausibility
logic, and abstention behavior.

\section{Overfitting, Plausibility, and Reward Hacking}

One of the most stable lessons from APR research is that passing the available tests is not the same as fixing
the underlying defect. Multiple studies in both classical APR and LLM-based APR report high rates of
overfitting, where a patch satisfies the visible oracle while remaining semantically wrong or behaviorally
fragile. More recent LLM repair studies continue to confirm the same pattern: some patches pass benchmark tests
yet fail stronger developer suites, and some even alter behavior in unintended ways while still looking
plausible under the benchmark evaluation protocol \cite{llmaprslr2024,aprtop2025,compass2026}.

For PatchOps, this literature is not background trivia. It directly motivates the decision to treat safety and
honesty as first-class research concerns. If the field already knows that plausible patches can be wrong, then a
CI repair system that automatically proposes or opens patches without explicit safeguards is not just
incomplete; it is methodologically naive. PatchOps responds to this by adding layers that go beyond basic patch
generation: reward-hack detection, patch-plausibility scoring, minimality-oriented gating, and abstention paths
for ambiguous cases.

This also helps explain why the thesis does not frame safety as a product add-on. Safety is not merely a user
experience preference. It is a scientific necessity forced by the known limitations of CI and APR evaluation
oracles. In this sense, the PatchOps safety gate is a direct descendant of the overfitting debate in APR.

\section{Benchmarks and Datasets}

The choice of benchmark shapes what claims can be made. For PatchOps, the most important benchmark is
\textbf{CI-Repair-Bench}, which offers 567 CI failure instances across more than one hundred Python repositories,
organized into twelve categories and validated through full workflow re-execution. Its main value is that it
turns PatchOps into a comparable system rather than an isolated prototype \cite{cirepairbench2026}. A thesis evaluated on
CI-Repair-Bench can be discussed alongside existing repair baselines, while a thesis evaluated only on a private
demo set would be much harder to defend.

However, no single dataset is enough for all PatchOps questions. CI-Repair-Bench is strong for end-to-end
repair success but weaker for the kinds of safety claims that matter most to this thesis. It does not richly
label risky conditions such as missing secrets, deployment permissions, or transient failures where rerunning is
better than patching. This is why PatchOps also relies conceptually on controlled safety and fault-injection
cases, as well as on broader CI log corpora such as \textbf{GHALogs}, which are useful for classification and
distributional context even if they do not directly support patch validation.

The literature also highlights the value of reproducible CI execution tooling, such as \textbf{GitBug-Actions},
which shows how GitHub Actions workflows can be replayed locally in controlled environments. This line of work
matters because practical CI evaluation is often bottlenecked not by the patch generator, but by the difficulty
of re-executing workflows repeatedly, cheaply, and reproducibly. PatchOps fits naturally into this methodological
space: its validation loop is meaningful only if the execution environment can be treated as dependable enough
for repeated evaluation \cite{ghalogs2025,gitbugactions2024}.

\section{Failure Taxonomy and the Diversity of CI Failures}

Another important strand of related work studies why GitHub Actions workflows fail in the first place.
Empirical taxonomy work, including the large-scale TOSEM 2025 study of GitHub Actions failures, shows that CI
failures are diverse and cannot be reduced to one or two failure modes. Project-level configuration errors,
dependency problems, runtime issues, environmental setup failures, permissions, and third-party service
interactions all appear in practice \cite{ghataxonomy2025}. This diversity is essential for PatchOps because it justifies both the
failure classifier and the decision policy.

If all failures were simply code-edit problems, then an always-patch system might be a defensible first
approximation. The literature shows the opposite. CI failures include many situations where the safest or most
appropriate action is not a patch. This is why PatchOps organizes failures into operational categories and why
the thesis places importance on taxonomy crosswalks between the project's own labels, CI-Repair-Bench's
categories, and broader empirical taxonomies from the literature.

This taxonomy perspective also strengthens the thesis gap. Public benchmarks emphasize repairable failures, but
the operational world includes many cases that are better treated as rerun, preview, or escalate. PatchOps
turns this mismatch into a contribution by explicitly modeling those decision classes rather than collapsing all
failures into the same repair workflow.

\section{Flaky and Transient Failures}

Research on flaky tests and transient CI failures provides one of the strongest motivations for PatchOps'
abstention-aware behavior. Studies on GitHub Actions flakiness show that some failures arise from instability in
external services, timing effects, environment differences, network issues, or order-dependent test behavior.
In such cases, producing a source-code patch is often the wrong response. A rerun, quarantine, or human review
is safer and more informative \cite{ghaflaky2026}.

This is a crucial distinction between PatchOps and simple repair pipelines. Most repair benchmarks reward only
successful patches; they do not reward the model for correctly deciding \emph{not} to patch. PatchOps instead
treats flaky detection and abstention as measurable successes. This creates a different notion of competence:
the system is not judged solely by how often it edits, but by how often it chooses the right operational action.

From a thesis perspective, this is one of the clearest novelty lanes available. Correctly distinguishing a
repairable failure from a transient one is both practically valuable and under-measured in prior CI repair
evaluations. It also produces compelling demos and easy-to-understand motivation, because a junk PR for a
network timeout is obviously less desirable than a correct rerun recommendation.

\section{Industry and Open-Source Landscape}

The broader industry landscape confirms that the product idea of ``AI fixes CI'' is no longer novel by itself.
GitHub Copilot Autofix, Dependabot, Renovate, Copilot coding-agent workflows, and multiple engineering blog
posts on self-healing CI all point in the same direction: software teams increasingly expect automated help when
pipelines fail. Some systems focus on vulnerabilities or dependency updates, some are issue-driven coding
agents, and others wire LLMs into branch-level repair loops \cite{logsage2025}.

This environment creates both a risk and an opportunity for PatchOps. The risk is that a vague thesis claim
such as ``we built an AI system that fixes CI'' will sound derivative and underspecified. The opportunity is
that most existing product narratives are not backed by public benchmarks, safety metrics, abstention analysis,
taxonomy grounding, or reproducible academic evaluation. In other words, the concept is being commoditized, but
the science is still underdeveloped.

PatchOps should therefore be understood as an academic response to an emerging product category. It is not
trying to beat commercial tools at deployment scale or enterprise integration. It is trying to provide the
measurement discipline, safety framing, and reproducibility that public product claims rarely supply.

\section{Implications for PatchOps Design}

Taken together, the literature suggests several design consequences for PatchOps.

First, log evidence extraction is necessary but insufficient. Systems such as LogSage show that careful log
processing improves diagnosis, but CI-Repair-Bench and later PatchOps results show that diagnosis alone does not
solve repair.

Second, repository-level reasoning is unavoidable. Real CI failures often require context beyond the visible log,
including workflow files, manifests, and source modules. This supports PatchOps' repository-aware design.

Third, validation must be embedded into the loop. The literature repeatedly shows that patch plausibility is a
weak signal of correctness. PatchOps therefore uses CI-style re-execution as part of its decision logic rather
than as an afterthought.

Fourth, abstention must be treated as a feature of competence, not as a failure of confidence. The taxonomy and
flakiness literature make it clear that some failures should not be patched automatically. PatchOps therefore
elevates rerun, preview, escalate, and abstain into first-class actions.

Fifth, safety and honesty checks are not optional. Overfitting and reward hacking are known risks in APR and
LLM repair. PatchOps responds by explicitly measuring unsafe actions, on-target behavior, and the difference
between apparently green and genuinely trustworthy outcomes.

\section{Chapter Conclusion and Positioning}

The literature reviewed in this chapter leads to a clear and disciplined positioning for PatchOps. Prior work
has already demonstrated that LLMs can assist with CI/CD diagnosis, that repository-aware CI repair can be
benchmarked through workflow re-execution, and that software repair pipelines benefit from structured rather
than purely free-form agent design. At the same time, prior work also makes clear that repair success remains
low, that test- or CI-passing patches can still be wrong, and that current benchmarks under-measure the
decision-theoretic aspects of safe automation.

PatchOps therefore does not claim novelty in patch generation alone. Its contribution lies in reframing CI
repair as a \emph{risk-sensitive, abstention-aware decision problem} built on top of repository context and
closed-loop validation. That framing is supported, not contradicted, by the literature. The next chapter turns
from this background to the formal problem definition, requirements, and scope of the PatchOps system.
